Heart failure remains a major cause of hospitalization, recurrent decompensation, and premature mortality worldwide, and it continues to impose a substantial clinical and public health burden \cite{khan2024global,bozkurt2025hfstats}. Prognosis in heart failure is highly heterogeneous, reflecting the combined influence of demographic characteristics, cardiac functional impairment, comorbidity, and laboratory abnormalities \cite{bozkurt2021universal,heidenreich2022guideline}. For this reason, accurate assessment of mortality risk is important in routine care, where it may inform treatment intensification, follow-up planning, monitoring priorities, and broader clinical decision-making \cite{heidenreich2022guideline,bozkurt2021universal}. Beyond estimating overall outcome probability, prognostic analysis is also valuable for clarifying which patient characteristics are most strongly associated with adverse events, thereby supporting a more structured understanding of disease severity and risk stratification \cite{khan2024global,bozkurt2025hfstats}.

Machine learning has been increasingly adopted in healthcare research for tasks such as diagnosis support, outcome prediction, and patient stratification \cite{alhejaily2024ai}. In comparison with conventional regression-based approaches, machine learning models can offer greater flexibility for capturing nonlinear relationships, threshold effects, and higher-order interactions among clinical variables \cite{saqib2024hfml,kokori2025survival}. Within heart failure research, these methods have been used to predict mortality, readmission, and other clinically relevant outcomes from routinely collected tabular data \cite{saqib2024hfml,kokori2025survival,liu2024hfhtn,ketabi2024mortality,chicco2020survival}. At the same time, predictive modeling alone is often insufficient in medical applications \cite{collins2024tripodai,moons2025probastai}. Clinically useful studies must also address interpretability by identifying the variables that consistently contribute to prognosis, ideally through complementary analytical perspectives rather than reliance on a single model-specific importance measure \cite{kundu2021explainable,caterson2024xai}.

Despite the growing literature, several methodological issues remain insufficiently resolved. Many studies report results for only a narrow subset of algorithms, limiting the ability to judge whether observed performance is robust to model choice \cite{xu2025meta,hajishah2025meta}. In addition, the effect of preprocessing decisions, including scaling and outlier handling, is often not examined systematically, even though such decisions may affect linear models, neural networks, and tree-based methods in different ways \cite{xu2025meta,collins2024tripodai}. A related challenge concerns the identification of influential risk factors in small clinical datasets, where limited sample size may lead to unstable estimates, overfitting, and model-dependent conclusions \cite{chicco2020survival,moons2025probastai}. Under these conditions, a more careful framework is needed that combines comparative prediction with multi-angle assessment of variable relevance using statistical testing, correlation structure, and model-based feature importance \cite{collins2024tripodai,kundu2021explainable,caterson2024xai}.

In this study, these issues are addressed through an integrated analysis of heart failure mortality prediction using logistic regression (LR), support vector machine (SVM), k-nearest neighbors (KNN), random forest (RF), extreme gradient boosting (XGBoost), and multilayer perceptron (MLP) models \cite{james2021islr,cortes1995svm,cover1967knn,breiman2001rf,chen2016xgboost,hastie2009esl}. Predictive performance was benchmarked across models, and preprocessing sensitivity was also examined to evaluate the effect of alternative data representations in a small-sample setting \cite{xu2025meta,hajishah2025meta}. In parallel, factors associated with \texttt{DEATH\_EVENT} were examined from multiple analytical perspectives, including descriptive comparison, univariate statistical testing, correlation analysis, and feature importance estimation, in order to identify clinically meaningful variables that showed consistent prognostic relevance. By combining model benchmarking with interpretable risk-factor analysis, this study aims to provide a methodologically transparent and clinically grounded assessment of mortality prediction in heart failure \cite{chicco2020survival,collins2024tripodai}.